%%%%%%
%
% $Autor: Wings $
% $Datum: 2020-01-18 11:15:45Z $
% $Pfad: WuSt/Skript/Produktspezifikation/powerpoint/ImageProcessing.tex $
% $Version: 4620 $
%
%%%%%%



\subsection{Starting up the Raspberry Pi Camera with the Jetson Nano}

The Raspberry Pi camera module is supported by the Jetson Nano Plug-and-Play, which makes commissioning easy. The device drivers are already installed with the operating system. Once the module is connected, the camera should show up as a device under \SHELL{/dev/video0}. If this is not the case, the plug connections must be checked. Especially during transport, the plug-in connection on the top of the board can come loose.

For the next steps the programme GStreamer is helpful \cite{GStreamer:2021}. The programme can be called as well as a programme, but it can also be integrated into your own programme. It enables the creation of single images and image sequences. GStreamer can be installed with the following commands:

\medskip

\SHELL{apt-get install libgstreamer1.0-0 gstreamer1.0-plugins-base   \textbackslash}

\SHELL{\quad gstreamer1.0-plugins-good gstreamer1.0-plugins-bad \textbackslash}

\SHELL{\quad gstreamer1.0-plugins-ugly gstreamer1.0-libav gstreamer1.0-doc \textbackslash}

\SHELL{\quad gstreamer1.0-tools gstreamer1.0-x gstreamer1.0-alsa \textbackslash}

\SHELL{\quad  gstreamer1.0-gl gstreamer1.0-gtk3 gstreamer1.0-qt5 \textbackslash}

\SHELL{\quad  gstreamer1.0-pulseaudio}

\medskip

Mit der Eingabe 

\medskip

\SHELL{\$ gst-launch-1.0 nvarguscamerasrc sensor\_id=0 ! nvoverlaysink}

\medskip

the camera can be tested. The command creates a camera image in full screen mode. To end the display, \keys{\ctrl + C} must be pressed.

\medskip

Other settings can also be made to the camera image. To request the GStreamer programme to open the live image of the camera with a width of 3820 pixels and a height of 2464 pixels at 21 frames per second and to display it in a window with a width of 960 pixels and a height of 616 pixels, the following line is entered:

\medskip

\SHELL{\$ gst-launch-1.0 nvarguscamerasrc sensor\_mode=0 !\textbackslash}

\SHELL{\quad 'video/x-raw(memory:NVMM),width=3820, height=2464,\textbackslash}

\SHELL{\quad framerate=21/1, format=NV12' ! nvvidconv flip-method=0 ! \textbackslash}

\SHELL{\quad 'video/x-raw,width=960, height=616' ! nvvidconv ! \textbackslash}
    
\SHELL{\quad nvegltransform ! nveglglessink -e}
%todo Apostrophe müssen gerade Striche sein, damit der Befehl per copy paste verwendet werden kann

When the camera is attached to the bracket on the Jetson Nano, the image may initially be upside down. If you set \SHELL{flip-method=2}, the image will be rotated by $180^\circ$. For more options to rotate and flip the image, see table~\ref{flip}.

\medskip

\medskip
\begin{table} [H]
    \centering
    \begin{tabular} {l c l}
        none                 & 0 & no rotation \\
        clockwise            & 1 & Rotate clockwise by $90^\circ$ \\
        rotate-180           & 2 & Rotate by $180^\circ$ \\
        counterclockwise     & 3 & Rotate counterclockwise by $90^\circ$ \\
        horizontal-flip      & 4 & Flip horizontally \\
        vertical-flip        & 5 & Flip vertical \\
        upper-left-diagonal  & 6 & flip over upper-left/lower-right diagonal \\
        upper-right-diagonal & 7 & Flip over upper right/lower left diagonal\\
        automatic            & 8 & flip method based on image-orientation tag\\
    \end{tabular}
    \caption{Operations to rotate and mirror the image in GStreamer}
    \label{Flip}
\end{table}


\medskip


The documentation for GSteamer can be found at \Cite{GStreamer:2021} and in a GitHub project. From the GitHub project, the documentation can be accessed via

\medskip

\SHELL{git clone https://gitlab.freedesktop.org/gstreamer/gst-docs}

\medskip

can be installed. Further plug-ins are also documented there.
